% --- Archon physics formalization source begin ---
% archon:physics
% archon:covers PhyXMiniProblems/problem_phyx_mini_0537.lean
% archon:source-report reports/phyx_mini/problem_phyx_mini_0537.source.json

\chapter{Physics problem phyx\_mini\_0537}
\label{ch:PhyXMiniProblems_problem_phyx_mini_0537}

\paragraph{Problem source.}
\#\# Physical scenario

A helium--neon laser can produce a green laser beam instead of a red one. Figure shows the transitions involved to form the red beam and the green beam. After a population inversion is established, neon atoms make a variety of downward transitions in falling from the state labeled \textbackslash{}( E\_4\textasciicircum{}* \textbackslash{}) down eventually to level \textbackslash{}( E\_1 \textbackslash{}) (arbitrarily assigned the energy \textbackslash{}( E\_1 = 0 \textbackslash{})). The atoms emit both red light with a wavelength of 632.8 nm in a transition \textbackslash{}( E\_4\textasciicircum{}* - E\_3 \textbackslash{}) and green light with a wavelength of 543 nm in a competing transition \textbackslash{}( E\_4\textasciicircum{}* - E\_2 \textbackslash{}).

\#\# Question

What is the energy \textbackslash{}( E\_2 \textbackslash{})? Assume the atoms are in a cavity between mirrors designed to reflect the green light with high efficiency but to allow the red light to leave the cavity immediately. Then stimulated emission can lead to the buildup of a collimated beam of green light between the mirrors having a greater intensity than that of the red light. To constitute the radiated laser beam, a small fraction of the green light is permitted to escape by transmission through one mirror. The mirrors forming the resonant cavity can be made of layers of silicon dioxide (index of refraction \textbackslash{}( n = 1.458 \textbackslash{})) and titanium dioxide (index of refraction varies between 1.9 and 2.6).

\#\# Answer choices

- A: A: \textbackslash{}(24.16 \textbackslash{}text\{ eV\}\textbackslash{})

- B: B: \textbackslash{}(22.01 \textbackslash{}text\{ eV\}\textbackslash{})

- C: C: \textbackslash{}(15.31 \textbackslash{}text\{ eV\}\textbackslash{})

- D: D: \textbackslash{}(18.37 \textbackslash{}text\{ eV\}\textbackslash{})

\#\# Figure caption (auxiliary; use the image as primary evidence)

The image is a diagram representing energy levels of an atom. It includes the following elements:

- Energy levels labeled \textbackslash{}(E\_1\textbackslash{}), \textbackslash{}(E\_2\textbackslash{}), \textbackslash{}(E\_3\textbackslash{}), and \textbackslash{}(E\_4\textasciicircum{}*\textbackslash{}).
- \textbackslash{}(E\_1\textbackslash{}) is marked as the "Ground state" with an energy of 0 eV.
- \textbackslash{}(E\_2\textbackslash{}) and \textbackslash{}(E\_3\textbackslash{}) are close together.
- \textbackslash{}(E\_3\textbackslash{}) is labeled with an energy of 18.70 eV.
- \textbackslash{}(E\_4\textasciicircum{}*\textbackslash{}) is labeled with an energy of 20.66 eV.
- There is a vertical arrow pointing upwards labeled "ENERGY."
- Two downward arrows from \textbackslash{}(E\_4\textasciicircum{}*\textbackslash{}) to \textbackslash{}(E\_3\textbackslash{}) and from \textbackslash{}(E\_3\textbackslash{}) to \textbackslash{}(E\_2\textbackslash{}) are depicted.
- The arrow from \textbackslash{}(E\_4\textasciicircum{}*\textbackslash{}) to \textbackslash{}(E\_3\textbackslash{}) is labeled "Red."
- The arrow from \textbackslash{}(E\_3\textbackslash{}) to \textbackslash{}(E\_2\textbackslash{}) is labeled "Green."

This diagram illustrates the transitions between different energy states along with the associated colors of emitted photons.

\#\# Dataset metadata

Quantum Phenomena; Physical Model Grounding Reasoning, Multi-Formula Reasoning

\paragraph{Recorded answer/context.}
D: D: \textbackslash{}(18.37 \textbackslash{}text\{ eV\}\textbackslash{})

\paragraph{Figure/image path.}
/root/proposal\_for\_physic/science-mango/phyx\_mini\_run/phyx\_data/test\_image/537.png

\paragraph{Formalization target.}
create a compiling Lean file with sorry bodies at `PhyXMiniProblems/problem\_phyx\_mini\_0537.lean`. The Lean declarations must preserve the physical quantities, dimensions or dimensional roles, figure labels, governing-law hypotheses, and final relation expressed by this problem.
Use Mathlib/Physlib names found through LeanExplore where available. If a physics API is missing, introduce faithful local abstractions rather than scalar placeholder aliases.

\begin{theorem}[Physics formalization target]
\label{thm:physics:phyx_mini_0537:target}
\lean{PhyXMiniProblems.ProblemPhyXMini0537.problem_phyx_mini_0537}
\uses{def:physics:phyx-mini-0537:phyxminiproblems-problemphyxmini0537-wavelengthinmeters, def:physics:phyx-mini-0537:phyxminiproblems-problemphyxmini0537-energyinjoules, def:physics:phyx-mini-0537:phyxminiproblems-problemphyxmini0537-energyinelectronvolts, def:physics:phyx-mini-0537:phyxminiproblems-problemphyxmini0537-planckconstantinjouleseconds, def:physics:phyx-mini-0537:phyxminiproblems-problemphyxmini0537-speedoflightinmeterspersecond, def:physics:phyx-mini-0537:phyxminiproblems-problemphyxmini0537-heliumneongreenlasersetup, def:physics:phyx-mini-0537:phyxminiproblems-problemphyxmini0537-matchesheliumneonlaserscenario, def:physics:phyx-mini-0537:phyxminiproblems-problemphyxmini0537-matchessuppliedenergylevelfigure, def:physics:phyx-mini-0537:phyxminiproblems-problemphyxmini0537-matchestransitionwavelengthreadouts, def:physics:phyx-mini-0537:phyxminiproblems-problemphyxmini0537-matchescavitydescription, def:physics:phyx-mini-0537:phyxminiproblems-problemphyxmini0537-hasphysicallaserquantities, def:physics:phyx-mini-0537:phyxminiproblems-problemphyxmini0537-usesordinaryplanckconstant, def:physics:phyx-mini-0537:phyxminiproblems-problemphyxmini0537-satisfiesradiativetransitionlaws, def:physics:phyx-mini-0537:phyxminiproblems-problemphyxmini0537-isuniquematchingdisplayedenergy}
The assigned autoformalize agent should translate this physics problem into Lean declarations in the covered file, with theorem and lemma proof bodies written as `by sorry`.
\end{theorem}
\begin{proof}
This is an autoformalization task, not a proof task. Produce statements that can later be proved by the physics prover without weakening the physical model.
\end{proof}

% --- Archon declaration topology begin ---
\section{Declaration topology}
\begin{definition}[energy Dimension]
  \label{def:physics:phyx-mini-0537:phyxminiproblems-problemphyxmini0537-energydimension}
  \lean{PhyXMiniProblems.ProblemPhyXMini0537.energyDimension}
  The physical dimension of energy, M L\textasciicircum{}2 T\textasciicircum{}-2
\end{definition}
\begin{proof}
  This is the declared model definition of energy Dimension.
\end{proof}
\begin{definition}[action Dimension]
  \label{def:physics:phyx-mini-0537:phyxminiproblems-problemphyxmini0537-actiondimension}
  \lean{PhyXMiniProblems.ProblemPhyXMini0537.actionDimension}
  \uses{def:physics:phyx-mini-0537:phyxminiproblems-problemphyxmini0537-energydimension}
  The physical dimension of action, energy multiplied by time
\end{definition}
\begin{proof}
  This is the declared model definition of action Dimension.
\end{proof}
\begin{definition}[Wavelength Quantity]
  \label{def:physics:phyx-mini-0537:phyxminiproblems-problemphyxmini0537-wavelengthquantity}
  \lean{PhyXMiniProblems.ProblemPhyXMini0537.WavelengthQuantity}
  A nonnegative, unit-independent physical wavelength
\end{definition}
\begin{proof}
  This is the declared model definition of Wavelength Quantity.
\end{proof}
\begin{definition}[Planck Constant Quantity]
  \label{def:physics:phyx-mini-0537:phyxminiproblems-problemphyxmini0537-planckconstantquantity}
  \lean{PhyXMiniProblems.ProblemPhyXMini0537.PlanckConstantQuantity}
  \uses{def:physics:phyx-mini-0537:phyxminiproblems-problemphyxmini0537-actiondimension}
  A nonnegative dimensionful quantity with the units of Planck's constant
\end{definition}
\begin{proof}
  This is the declared model definition of Planck Constant Quantity.
\end{proof}
\begin{definition}[Optical Intensity Quantity]
  \label{def:physics:phyx-mini-0537:phyxminiproblems-problemphyxmini0537-opticalintensityquantity}
  \lean{PhyXMiniProblems.ProblemPhyXMini0537.OpticalIntensityQuantity}
  Optical intensity, with physical dimension M T\textasciicircum{}-3
\end{definition}
\begin{proof}
  This is the declared model definition of Optical Intensity Quantity.
\end{proof}
\begin{definition}[wavelength Readout]
  \label{def:physics:phyx-mini-0537:phyxminiproblems-problemphyxmini0537-wavelengthreadout}
  \lean{PhyXMiniProblems.ProblemPhyXMini0537.wavelengthReadout}
  \uses{def:physics:phyx-mini-0537:phyxminiproblems-problemphyxmini0537-wavelengthquantity}
  Read a physical wavelength in a selected length unit
\end{definition}
\begin{proof}
  This is the declared model definition of wavelength Readout.
\end{proof}
\begin{definition}[wavelength In Meters]
  \label{def:physics:phyx-mini-0537:phyxminiproblems-problemphyxmini0537-wavelengthinmeters}
  \lean{PhyXMiniProblems.ProblemPhyXMini0537.wavelengthInMeters}
  \uses{def:physics:phyx-mini-0537:phyxminiproblems-problemphyxmini0537-wavelengthquantity, def:physics:phyx-mini-0537:phyxminiproblems-problemphyxmini0537-wavelengthreadout}
  Read a wavelength in coherent-SI metres
\end{definition}
\begin{proof}
  This is the declared model definition of wavelength In Meters.
\end{proof}
\begin{definition}[wavelength In Nanometers]
  \label{def:physics:phyx-mini-0537:phyxminiproblems-problemphyxmini0537-wavelengthinnanometers}
  \lean{PhyXMiniProblems.ProblemPhyXMini0537.wavelengthInNanometers}
  \uses{def:physics:phyx-mini-0537:phyxminiproblems-problemphyxmini0537-wavelengthquantity, def:physics:phyx-mini-0537:phyxminiproblems-problemphyxmini0537-wavelengthreadout}
  Read a wavelength in nanometres
\end{definition}
\begin{proof}
  This is the declared model definition of wavelength In Nanometers.
\end{proof}
\begin{definition}[energy In Joules]
  \label{def:physics:phyx-mini-0537:phyxminiproblems-problemphyxmini0537-energyinjoules}
  \lean{PhyXMiniProblems.ProblemPhyXMini0537.energyInJoules}
  Read a physical energy in coherent-SI joules
\end{definition}
\begin{proof}
  This is the declared model definition of energy In Joules.
\end{proof}
\begin{definition}[energy In Electron Volts]
  \label{def:physics:phyx-mini-0537:phyxminiproblems-problemphyxmini0537-energyinelectronvolts}
  \lean{PhyXMiniProblems.ProblemPhyXMini0537.energyInElectronVolts}
  \uses{def:physics:phyx-mini-0537:phyxminiproblems-problemphyxmini0537-energyinjoules}
  Read an energy in electron volts using Physlib's electron-volt quantity
\end{definition}
\begin{proof}
  This is the declared model definition of energy In Electron Volts.
\end{proof}
\begin{definition}[planck Constant In Joule Seconds]
  \label{def:physics:phyx-mini-0537:phyxminiproblems-problemphyxmini0537-planckconstantinjouleseconds}
  \lean{PhyXMiniProblems.ProblemPhyXMini0537.planckConstantInJouleSeconds}
  \uses{def:physics:phyx-mini-0537:phyxminiproblems-problemphyxmini0537-planckconstantquantity}
  Read a dimensionful Planck constant in coherent-SI joule-seconds
\end{definition}
\begin{proof}
  This is the declared model definition of planck Constant In Joule Seconds.
\end{proof}
\begin{definition}[speed Of Light In Meters Per Second]
  \label{def:physics:phyx-mini-0537:phyxminiproblems-problemphyxmini0537-speedoflightinmeterspersecond}
  \lean{PhyXMiniProblems.ProblemPhyXMini0537.speedOfLightInMetersPerSecond}
  Physlib's dimensionful speed of light, read in SI metres per second
\end{definition}
\begin{proof}
  This is the declared model definition of speed Of Light In Meters Per Second.
\end{proof}
\begin{definition}[optical Intensity In SI]
  \label{def:physics:phyx-mini-0537:phyxminiproblems-problemphyxmini0537-opticalintensityinsi}
  \lean{PhyXMiniProblems.ProblemPhyXMini0537.opticalIntensityInSI}
  \uses{def:physics:phyx-mini-0537:phyxminiproblems-problemphyxmini0537-opticalintensityquantity}
  Read a physical optical intensity in coherent SI units
\end{definition}
\begin{proof}
  This is the declared model definition of optical Intensity In SI.
\end{proof}
\begin{definition}[ordinary Planck Constant SIValue]
  \label{def:physics:phyx-mini-0537:phyxminiproblems-problemphyxmini0537-ordinaryplanckconstantsivalue}
  \lean{PhyXMiniProblems.ProblemPhyXMini0537.ordinaryPlanckConstantSIValue}
  The exact coherent-SI calibration of the ordinary Planck constant
\end{definition}
\begin{proof}
  This is the declared model definition of ordinary Planck Constant SIValue.
\end{proof}
\begin{definition}[energy Level Display Tolerance In Electron Volts]
  \label{def:physics:phyx-mini-0537:phyxminiproblems-problemphyxmini0537-energyleveldisplaytoleranceinelectronvolts}
  \lean{PhyXMiniProblems.ProblemPhyXMini0537.energyLevelDisplayToleranceInElectronVolts}
  The displayed level energies have two digits after the decimal point, the red wavelength has one, and the green wavelength is printed to the nearest nanometre. These half-unit tolerances represent the precision of those labels
\end{definition}
\begin{proof}
  This is the declared model definition of energy Level Display Tolerance In Electron Volts.
\end{proof}
\begin{definition}[red Wavelength Display Tolerance In Nanometers]
  \label{def:physics:phyx-mini-0537:phyxminiproblems-problemphyxmini0537-redwavelengthdisplaytoleranceinnanometers}
  \lean{PhyXMiniProblems.ProblemPhyXMini0537.redWavelengthDisplayToleranceInNanometers}
  This def records the red Wavelength Display Tolerance In Nanometers component of the problem-specific physical model
\end{definition}
\begin{proof}
  This is the declared model definition of red Wavelength Display Tolerance In Nanometers.
\end{proof}
\begin{definition}[green Wavelength Display Tolerance In Nanometers]
  \label{def:physics:phyx-mini-0537:phyxminiproblems-problemphyxmini0537-greenwavelengthdisplaytoleranceinnanometers}
  \lean{PhyXMiniProblems.ProblemPhyXMini0537.greenWavelengthDisplayToleranceInNanometers}
  This def records the green Wavelength Display Tolerance In Nanometers component of the problem-specific physical model
\end{definition}
\begin{proof}
  This is the declared model definition of green Wavelength Display Tolerance In Nanometers.
\end{proof}
\begin{definition}[Atomic Energy Level]
  \label{def:physics:phyx-mini-0537:phyxminiproblems-problemphyxmini0537-atomicenergylevel}
  \lean{PhyXMiniProblems.ProblemPhyXMini0537.AtomicEnergyLevel}
  The four horizontal neon energy levels named in the figure
\end{definition}
\begin{proof}
  This is the declared model definition of Atomic Energy Level.
\end{proof}
\begin{definition}[Emission Transition]
  \label{def:physics:phyx-mini-0537:phyxminiproblems-problemphyxmini0537-emissiontransition}
  \lean{PhyXMiniProblems.ProblemPhyXMini0537.EmissionTransition}
  The two competing downward transitions described by the problem
\end{definition}
\begin{proof}
  This is the declared model definition of Emission Transition.
\end{proof}
\begin{definition}[Laser Color]
  \label{def:physics:phyx-mini-0537:phyxminiproblems-problemphyxmini0537-lasercolor}
  \lean{PhyXMiniProblems.ProblemPhyXMini0537.LaserColor}
  The color label printed beside an emission arrow
\end{definition}
\begin{proof}
  This is the declared model definition of Laser Color.
\end{proof}
\begin{definition}[transition Initial Level]
  \label{def:physics:phyx-mini-0537:phyxminiproblems-problemphyxmini0537-transitioninitiallevel}
  \lean{PhyXMiniProblems.ProblemPhyXMini0537.transitionInitialLevel}
  \uses{def:physics:phyx-mini-0537:phyxminiproblems-problemphyxmini0537-atomicenergylevel, def:physics:phyx-mini-0537:phyxminiproblems-problemphyxmini0537-emissiontransition}
  Both depicted emission transitions start from the populated E4* level
\end{definition}
\begin{proof}
  This is the declared model definition of transition Initial Level.
\end{proof}
\begin{definition}[transition Final Level]
  \label{def:physics:phyx-mini-0537:phyxminiproblems-problemphyxmini0537-transitionfinallevel}
  \lean{PhyXMiniProblems.ProblemPhyXMini0537.transitionFinalLevel}
  \uses{def:physics:phyx-mini-0537:phyxminiproblems-problemphyxmini0537-atomicenergylevel, def:physics:phyx-mini-0537:phyxminiproblems-problemphyxmini0537-emissiontransition}
  The red transition ends at E3; the green transition ends at E2
\end{definition}
\begin{proof}
  This is the declared model definition of transition Final Level.
\end{proof}
\begin{definition}[transition Color]
  \label{def:physics:phyx-mini-0537:phyxminiproblems-problemphyxmini0537-transitioncolor}
  \lean{PhyXMiniProblems.ProblemPhyXMini0537.transitionColor}
  \uses{def:physics:phyx-mini-0537:phyxminiproblems-problemphyxmini0537-emissiontransition, def:physics:phyx-mini-0537:phyxminiproblems-problemphyxmini0537-lasercolor}
  Color associated with each of the two transitions
\end{definition}
\begin{proof}
  This is the declared model definition of transition Color.
\end{proof}
\begin{definition}[Neon Energy Level Figure]
  \label{def:physics:phyx-mini-0537:phyxminiproblems-problemphyxmini0537-neonenergylevelfigure}
  \lean{PhyXMiniProblems.ProblemPhyXMini0537.NeonEnergyLevelFigure}
  \uses{def:physics:phyx-mini-0537:phyxminiproblems-problemphyxmini0537-atomicenergylevel, def:physics:phyx-mini-0537:phyxminiproblems-problemphyxmini0537-emissiontransition, def:physics:phyx-mini-0537:phyxminiproblems-problemphyxmini0537-lasercolor}
  Presentation data visible in image 537. An optional printed energy records that the raster deliberately does not print the requested energy of E2
\end{definition}
\begin{proof}
  This is the declared model definition of Neon Energy Level Figure.
\end{proof}
\begin{definition}[Laser Medium Kind]
  \label{def:physics:phyx-mini-0537:phyxminiproblems-problemphyxmini0537-lasermediumkind}
  \lean{PhyXMiniProblems.ProblemPhyXMini0537.LaserMediumKind}
  The gain medium named in the problem statement
\end{definition}
\begin{proof}
  This is the declared model definition of Laser Medium Kind.
\end{proof}
\begin{definition}[Mirror Layer Material]
  \label{def:physics:phyx-mini-0537:phyxminiproblems-problemphyxmini0537-mirrorlayermaterial}
  \lean{PhyXMiniProblems.ProblemPhyXMini0537.MirrorLayerMaterial}
  Dielectric materials named for the multilayer cavity mirrors
\end{definition}
\begin{proof}
  This is the declared model definition of Mirror Layer Material.
\end{proof}
\begin{definition}[Resonant Cavity Data]
  \label{def:physics:phyx-mini-0537:phyxminiproblems-problemphyxmini0537-resonantcavitydata}
  \lean{PhyXMiniProblems.ProblemPhyXMini0537.ResonantCavityData}
  \uses{def:physics:phyx-mini-0537:phyxminiproblems-problemphyxmini0537-lasercolor, def:physics:phyx-mini-0537:phyxminiproblems-problemphyxmini0537-mirrorlayermaterial}
  Material and optical data for the resonant cavity. Reflectivity, transmittance, and refractive index are dimensionless. Boolean fields retain qualitative phrases for which the source gives no numerical threshold
\end{definition}
\begin{proof}
  This is the declared model definition of Resonant Cavity Data.
\end{proof}
\begin{definition}[Helium Neon Green Laser Setup]
  \label{def:physics:phyx-mini-0537:phyxminiproblems-problemphyxmini0537-heliumneongreenlasersetup}
  \lean{PhyXMiniProblems.ProblemPhyXMini0537.HeliumNeonGreenLaserSetup}
  \uses{def:physics:phyx-mini-0537:phyxminiproblems-problemphyxmini0537-wavelengthquantity, def:physics:phyx-mini-0537:phyxminiproblems-problemphyxmini0537-planckconstantquantity, def:physics:phyx-mini-0537:phyxminiproblems-problemphyxmini0537-opticalintensityquantity, def:physics:phyx-mini-0537:phyxminiproblems-problemphyxmini0537-atomicenergylevel, def:physics:phyx-mini-0537:phyxminiproblems-problemphyxmini0537-emissiontransition, def:physics:phyx-mini-0537:phyxminiproblems-problemphyxmini0537-lasercolor, def:physics:phyx-mini-0537:phyxminiproblems-problemphyxmini0537-neonenergylevelfigure, def:physics:phyx-mini-0537:phyxminiproblems-problemphyxmini0537-lasermediumkind, def:physics:phyx-mini-0537:phyxminiproblems-problemphyxmini0537-resonantcavitydata}
  Independent physical data for the laser. In particular, energyAtLevel .e2 has no fixed value in this structure
\end{definition}
\begin{proof}
  This is the declared model definition of Helium Neon Green Laser Setup.
\end{proof}
\begin{definition}[Matches Helium Neon Laser Scenario]
  \label{def:physics:phyx-mini-0537:phyxminiproblems-problemphyxmini0537-matchesheliumneonlaserscenario}
  \lean{PhyXMiniProblems.ProblemPhyXMini0537.MatchesHeliumNeonLaserScenario}
  \uses{def:physics:phyx-mini-0537:phyxminiproblems-problemphyxmini0537-opticalintensityinsi, def:physics:phyx-mini-0537:phyxminiproblems-problemphyxmini0537-heliumneongreenlasersetup}
  The qualitative atomic and laser behavior stated in the prose
\end{definition}
\begin{proof}
  This is the declared model definition of Matches Helium Neon Laser Scenario.
\end{proof}
\begin{definition}[Matches Supplied Energy Level Figure]
  \label{def:physics:phyx-mini-0537:phyxminiproblems-problemphyxmini0537-matchessuppliedenergylevelfigure}
  \lean{PhyXMiniProblems.ProblemPhyXMini0537.MatchesSuppliedEnergyLevelFigure}
  \uses{def:physics:phyx-mini-0537:phyxminiproblems-problemphyxmini0537-energyinelectronvolts, def:physics:phyx-mini-0537:phyxminiproblems-problemphyxmini0537-energyleveldisplaytoleranceinelectronvolts, def:physics:phyx-mini-0537:phyxminiproblems-problemphyxmini0537-transitioninitiallevel, def:physics:phyx-mini-0537:phyxminiproblems-problemphyxmini0537-transitionfinallevel, def:physics:phyx-mini-0537:phyxminiproblems-problemphyxmini0537-transitioncolor, def:physics:phyx-mini-0537:phyxminiproblems-problemphyxmini0537-heliumneongreenlasersetup}
  Primary-raster evidence. The arrow endpoints record that green is E4* -> E2, as shown in the image, rather than E3 -> E2 from the inconsistent auxiliary caption. The unknown E2 has no printed numerical label
\end{definition}
\begin{proof}
  This is the declared model definition of Matches Supplied Energy Level Figure.
\end{proof}
\begin{definition}[Matches Transition Wavelength Readouts]
  \label{def:physics:phyx-mini-0537:phyxminiproblems-problemphyxmini0537-matchestransitionwavelengthreadouts}
  \lean{PhyXMiniProblems.ProblemPhyXMini0537.MatchesTransitionWavelengthReadouts}
  \uses{def:physics:phyx-mini-0537:phyxminiproblems-problemphyxmini0537-wavelengthinnanometers, def:physics:phyx-mini-0537:phyxminiproblems-problemphyxmini0537-redwavelengthdisplaytoleranceinnanometers, def:physics:phyx-mini-0537:phyxminiproblems-problemphyxmini0537-greenwavelengthdisplaytoleranceinnanometers, def:physics:phyx-mini-0537:phyxminiproblems-problemphyxmini0537-heliumneongreenlasersetup}
  The two vacuum wavelengths stated in the problem prose
\end{definition}
\begin{proof}
  This is the declared model definition of Matches Transition Wavelength Readouts.
\end{proof}
\begin{definition}[Matches Cavity Description]
  \label{def:physics:phyx-mini-0537:phyxminiproblems-problemphyxmini0537-matchescavitydescription}
  \lean{PhyXMiniProblems.ProblemPhyXMini0537.MatchesCavityDescription}
  \uses{def:physics:phyx-mini-0537:phyxminiproblems-problemphyxmini0537-heliumneongreenlasersetup}
  The stated dielectric materials and their qualitative cavity behavior
\end{definition}
\begin{proof}
  This is the declared model definition of Matches Cavity Description.
\end{proof}
\begin{definition}[Has Physical Laser Quantities]
  \label{def:physics:phyx-mini-0537:phyxminiproblems-problemphyxmini0537-hasphysicallaserquantities}
  \lean{PhyXMiniProblems.ProblemPhyXMini0537.HasPhysicalLaserQuantities}
  \uses{def:physics:phyx-mini-0537:phyxminiproblems-problemphyxmini0537-wavelengthinmeters, def:physics:phyx-mini-0537:phyxminiproblems-problemphyxmini0537-energyinjoules, def:physics:phyx-mini-0537:phyxminiproblems-problemphyxmini0537-planckconstantinjouleseconds, def:physics:phyx-mini-0537:phyxminiproblems-problemphyxmini0537-heliumneongreenlasersetup}
  Positivity and ordering conditions for the physical quantities
\end{definition}
\begin{proof}
  This is the declared model definition of Has Physical Laser Quantities.
\end{proof}
\begin{definition}[Uses Ordinary Planck Constant]
  \label{def:physics:phyx-mini-0537:phyxminiproblems-problemphyxmini0537-usesordinaryplanckconstant}
  \lean{PhyXMiniProblems.ProblemPhyXMini0537.UsesOrdinaryPlanckConstant}
  \uses{def:physics:phyx-mini-0537:phyxminiproblems-problemphyxmini0537-planckconstantinjouleseconds, def:physics:phyx-mini-0537:phyxminiproblems-problemphyxmini0537-ordinaryplanckconstantsivalue, def:physics:phyx-mini-0537:phyxminiproblems-problemphyxmini0537-heliumneongreenlasersetup}
  The ordinary Planck constant with its standard coherent-SI calibration
\end{definition}
\begin{proof}
  This is the declared model definition of Uses Ordinary Planck Constant.
\end{proof}
\begin{definition}[Satisfies Radiative Transition Laws]
  \label{def:physics:phyx-mini-0537:phyxminiproblems-problemphyxmini0537-satisfiesradiativetransitionlaws}
  \lean{PhyXMiniProblems.ProblemPhyXMini0537.SatisfiesRadiativeTransitionLaws}
  \uses{def:physics:phyx-mini-0537:phyxminiproblems-problemphyxmini0537-wavelengthinmeters, def:physics:phyx-mini-0537:phyxminiproblems-problemphyxmini0537-energyinjoules, def:physics:phyx-mini-0537:phyxminiproblems-problemphyxmini0537-planckconstantinjouleseconds, def:physics:phyx-mini-0537:phyxminiproblems-problemphyxmini0537-speedoflightinmeterspersecond, def:physics:phyx-mini-0537:phyxminiproblems-problemphyxmini0537-transitioninitiallevel, def:physics:phyx-mini-0537:phyxminiproblems-problemphyxmini0537-transitionfinallevel, def:physics:phyx-mini-0537:phyxminiproblems-problemphyxmini0537-heliumneongreenlasersetup}
  Uniform laws for every allowed transition: the photon carries the atomic energy drop and obeys E\_gamma = h c / lambda. Neither law singles out E2 or contains an answer-choice value
\end{definition}
\begin{proof}
  This is the declared model definition of Satisfies Radiative Transition Laws.
\end{proof}
\begin{definition}[Answer Choice]
  \label{def:physics:phyx-mini-0537:phyxminiproblems-problemphyxmini0537-answerchoice}
  \lean{PhyXMiniProblems.ProblemPhyXMini0537.AnswerChoice}
  Labels of the four energy choices supplied with the problem
\end{definition}
\begin{proof}
  This is the declared model definition of Answer Choice.
\end{proof}
\begin{definition}[displayed Energy In Electron Volts]
  \label{def:physics:phyx-mini-0537:phyxminiproblems-problemphyxmini0537-displayedenergyinelectronvolts}
  \lean{PhyXMiniProblems.ProblemPhyXMini0537.displayedEnergyInElectronVolts}
  \uses{def:physics:phyx-mini-0537:phyxminiproblems-problemphyxmini0537-answerchoice}
  Energy in electron volts printed beside each answer label
\end{definition}
\begin{proof}
  This is the declared model definition of displayed Energy In Electron Volts.
\end{proof}
\begin{definition}[recorded Dataset Answer]
  \label{def:physics:phyx-mini-0537:phyxminiproblems-problemphyxmini0537-recordeddatasetanswer}
  \lean{PhyXMiniProblems.ProblemPhyXMini0537.recordedDatasetAnswer}
  \uses{def:physics:phyx-mini-0537:phyxminiproblems-problemphyxmini0537-answerchoice}
  The answer label recorded in the source dataset, retained as metadata
\end{definition}
\begin{proof}
  This is the declared model definition of recorded Dataset Answer.
\end{proof}
\begin{definition}[displayed Energy Compatibility Tolerance]
  \label{def:physics:phyx-mini-0537:phyxminiproblems-problemphyxmini0537-displayedenergycompatibilitytolerance}
  \lean{PhyXMiniProblems.ProblemPhyXMini0537.displayedEnergyCompatibilityTolerance}
  The rounded 20.66 eV and 543 nm readouts propagate to less than two hundredths of an electron volt in the derived E2 value
\end{definition}
\begin{proof}
  This is the declared model definition of displayed Energy Compatibility Tolerance.
\end{proof}
\begin{definition}[Matches Displayed Energy]
  \label{def:physics:phyx-mini-0537:phyxminiproblems-problemphyxmini0537-matchesdisplayedenergy}
  \lean{PhyXMiniProblems.ProblemPhyXMini0537.MatchesDisplayedEnergy}
  \uses{def:physics:phyx-mini-0537:phyxminiproblems-problemphyxmini0537-energyinelectronvolts, def:physics:phyx-mini-0537:phyxminiproblems-problemphyxmini0537-heliumneongreenlasersetup, def:physics:phyx-mini-0537:phyxminiproblems-problemphyxmini0537-answerchoice, def:physics:phyx-mini-0537:phyxminiproblems-problemphyxmini0537-displayedenergyinelectronvolts, def:physics:phyx-mini-0537:phyxminiproblems-problemphyxmini0537-displayedenergycompatibilitytolerance}
  The physical E2 energy is compatible with a printed energy choice
\end{definition}
\begin{proof}
  This is the declared model definition of Matches Displayed Energy.
\end{proof}
\begin{definition}[Is Unique Matching Displayed Energy]
  \label{def:physics:phyx-mini-0537:phyxminiproblems-problemphyxmini0537-isuniquematchingdisplayedenergy}
  \lean{PhyXMiniProblems.ProblemPhyXMini0537.IsUniqueMatchingDisplayedEnergy}
  \uses{def:physics:phyx-mini-0537:phyxminiproblems-problemphyxmini0537-heliumneongreenlasersetup, def:physics:phyx-mini-0537:phyxminiproblems-problemphyxmini0537-answerchoice, def:physics:phyx-mini-0537:phyxminiproblems-problemphyxmini0537-matchesdisplayedenergy}
  A choice is the unique displayed energy compatible with physical E2
\end{definition}
\begin{proof}
  This is the declared model definition of Is Unique Matching Displayed Energy.
\end{proof}
\begin{lemma}[e2 energy In Joules exact]
  \label{lem:physics:phyx-mini-0537:phyxminiproblems-problemphyxmini0537-e2-energyinjoules-exact}
  \lean{PhyXMiniProblems.ProblemPhyXMini0537.e2_energyInJoules_exact}
  \uses{def:physics:phyx-mini-0537:phyxminiproblems-problemphyxmini0537-wavelengthinmeters, def:physics:phyx-mini-0537:phyxminiproblems-problemphyxmini0537-energyinjoules, def:physics:phyx-mini-0537:phyxminiproblems-problemphyxmini0537-planckconstantinjouleseconds, def:physics:phyx-mini-0537:phyxminiproblems-problemphyxmini0537-speedoflightinmeterspersecond, def:physics:phyx-mini-0537:phyxminiproblems-problemphyxmini0537-heliumneongreenlasersetup, def:physics:phyx-mini-0537:phyxminiproblems-problemphyxmini0537-matchesheliumneonlaserscenario, def:physics:phyx-mini-0537:phyxminiproblems-problemphyxmini0537-satisfiesradiativetransitionlaws}
  The green transition gives the symbolic value of E2 by subtracting the green photon's h c / lambda energy from the upper level
\end{lemma}
\begin{proof}
  The conclusion follows from the governing relations and figure readouts named in the statement of e2 energy In Joules exact.
\end{proof}
% --- Archon declaration topology end ---
% --- Archon physics formalization source end ---
